\documentclass[crop=false]{standalone}
\usepackage{15150toolbox}
\usepackage{comment}
\usepackage{needspace}

% Toggle: set \soltrue to render solutions inline (blue boxes).
\newif\ifsol
\soltrue   % comment this line for the student version

% Solution environment.
\ifsol
  \newenvironment{solution}
    {\par\color{solution_color}\begin{framed}\textbf{Solution.}\par}
    {\end{framed}}
\else
  \excludecomment{solution}
\fi

% Lecture-slide macros not in 15150toolbox.
\newcommand{\term}[1]{\textbf{\textit{#1}}}
\newcommand{\thmBox}[2]{\par\noindent\begin{mdframed}[backgroundcolor=blue!8,hidealllines=true]\textbf{Theorem#1.} #2\end{mdframed}}
\newcommand{\bcBox}[1]{\par\noindent\textbf{Base case.} #1\par}
\newcommand{\ihBox}[2]{\par\noindent\textbf{Inductive hypothesis#1.} #2\par}
\newcommand{\isBox}[2]{\par\noindent\textbf{Inductive step#1.} #2\par}

% Problem header.
\newcommand{\quizproblem}[3]{%
  \needspace{5\baselineskip}%
  \stepcounter{problem}%
  \ifsol
    \subsection*{Problem \arabic{problem}. #1 \hfill \timing{#2} \quad (#3 pts)}%
  \else
    \subsection*{Problem \arabic{problem}. #1 \hfill (#3 pts)}%
  \fi
}

\newcommand{\subpts}[1]{\textbf{(#1 pts)}}

\newcommand{\answerspace}[1]{\ifsol\else\vspace{#1}\par\fi}

\printout{Quiz 6}

\begin{document}

\begin{center}
  {\Large \textbf{15-150 --- Quiz 6}} \\
  \vspace{4pt}
  Topics: \textit{Sequences}; \textit{Lazy Programming};
  \textit{Imperative Programming}; \textit{Asynchronous Programming} \\
  \vspace{4pt}
  \textit{60 minutes.}
\end{center}

\vspace{1em}

\noindent\textbf{Name:} \hrulefill \hspace{2em} \textbf{Andrew ID:} \hrulefill

\vspace{1em}

% =============================================================================
% PROBLEM 1 --- Sequences: cost of a list-style algorithm on a sequence
% =============================================================================

\quizproblem{Cons and Effect}{8}{5}

Consider the following function on sequences:

\begin{codeblock}
  fun rev (S : 'a seq) : 'a seq =
    case Seq.length S of
      0 => Seq.empty ()
    | n => Seq.cons (Seq.nth S (n - 1)) (rev (Seq.take S (n - 1)))
\end{codeblock}

On lists this pattern is $O(n)$ work. On sequences it is not. Recall that
\code{Seq.cons} and \code{Seq.take} are each $O(n)$ work (both copy the
whole sequence), while \code{Seq.nth} and \code{Seq.length} are $O(1)$.
Let \code{n = Seq.length S}.

\textbf{(a)} \subpts{2} Give a tight big-$O$ bound for the total
\textbf{work} of \code{rev S}, and justify it in one sentence.

\answerspace{4cm}

\begin{solution}
\textbf{Work: $O(n^2)$.} There are $n$ recursive calls (the sequence
shrinks by one each time), and each does a \code{Seq.cons} and a
\code{Seq.take}, both $O(n)$. That is $O(n)$ work per call over $n$
calls: $O(n^2)$.

\textbf{Rubric:} 1, $O(n^2)$; 1, $n$ calls $\times$ $O(n)$ each.
\end{solution}

\textbf{(b)} \subpts{3} Implement \code{rev} a better way --- without
using \code{Seq.rev}. You should have an improved bound.

\answerspace{4cm}

\begin{solution}
\begin{codeblock}
  fun rev (S : 'a seq) : 'a seq =
    Seq.tabulate (fn i => Seq.nth S (Seq.length S - 1 - i))
                 (Seq.length S)
\end{codeblock}
Work is $O(n)$: a single \code{tabulate} with an $O(1)$ body, since
\code{nth} and \code{length} are constant time. (Span is $O(1)$, as the
elements are independent.)

\textbf{Rubric:} 2, \code{tabulate} over \code{n} with body
\code{nth S (n - 1 - i)}; 1, work $O(n)$.
\end{solution}

\newpage

% =============================================================================
% PROBLEM 2 --- Lazy: stream filter productivity
% =============================================================================

\quizproblem{Filter Failure}{6}{4}

Recall streams and the \code{filter} function on them:

\begin{codeblock}
  datatype 'a stream = Stream of (unit -> 'a front)
       and 'a front  = Empty | Cons of 'a * 'a stream

  fun filter p S =
    delay (fn () =>
      case expose S of
        Empty        => Empty
      | Cons (x, S') => if p x then Cons (x, filter p S')
                        else expose (filter p S'))
\end{codeblock}

A stream is \term{productive} if exposing it always terminates (returning
\code{Empty} or a \code{Cons}). Let \code{nats} be the productive infinite
stream \code{0, 1, 2, 3, ...}.

\textbf{(a)} \subpts{2} Consider \code{filter (fn x => x < 0) nats}.
Exposing this stream does not terminate, even though \code{nats} is
productive and \code{(fn x => x < 0)} is total. In one sentence, say why.

\answerspace{3cm}

\begin{solution}
No natural number is negative, so every element takes the \code{else}
branch: each expose just exposes the next recursive \code{filter} call,
searching for a match that never appears, so it never returns a front.

\textbf{Rubric:} 2, no element ever satisfies \code{p}, so the search
never terminates.
\end{solution}

\textbf{(b)} \subpts{2} Is \code{filter (fn x => x mod 2 = 0) nats}
productive? Answer yes or no, and justify in one sentence.

\answerspace{2.5cm}

\begin{solution}
\textbf{Yes.} Even numbers occur infinitely often in \code{nats}, so
every search in the \code{else} branch finds the next even number in
finitely many steps, and every expose terminates.

\textbf{Rubric:} 1, yes; 1, evens occur infinitely often, so each search
terminates.
\end{solution}

\newpage

% =============================================================================
% PROBLEM 3 --- Imperative: wrong number of ref cells (repeat)
% =============================================================================

\quizproblem{Store Wars}{6}{6}

A student writes \code{repeat (s, n)}, intended to build the string
\code{s} concatenated \code{n} times, by appending into a reference
cell. Knowing that refs shouldn't be global, they allocate it inside the
function:

\begin{codeblock}
  fun repeat (s, n) =
    let
      val buf = ref ""
    in
      if n = 0 then !buf
      else ( buf := !buf ^ s; repeat (s, n - 1) )
    end
\end{codeblock}

\textbf{(a)} \subpts{4} What does \code{repeat ("ab", 2)} evaluate to?
Explain in one or two sentences.

\answerspace{4cm}

\begin{solution}
\textbf{\code{""}}, for any arguments. Every call to \code{repeat},
including each recursive call, runs \code{ref ""} and allocates its own
cell: one call to \code{ref}, one box. Each \code{:=} writes into a cell
the next call abandons, and the base case reads a fresh empty one.

\textbf{Rubric:} 2, \code{""}; 2, each recursive call allocates its own
\code{buf}, losing the appends.
\end{solution}

\textbf{(b)} \subpts{2} In one sentence (or with code), how would you fix
\code{repeat} --- still using a ref --- so that it returns
\code{"abab"}?

\answerspace{3cm}

\begin{solution}
Allocate the ref \emph{once per top-level call}, and do the recursion in
an inner helper that shares that one cell:

\begin{codeblock}
  fun repeat (s, n) =
    let
      val buf = ref ""
      fun go 0 = !buf
        | go n = ( buf := !buf ^ s; go (n - 1) )
    in
      go n
    end
\end{codeblock}

\textbf{Rubric:} 2, one \code{ref} outside the recursion, inner helper
shares it (sentence or code).
\end{solution}

\newpage

% =============================================================================
% PROBLEM 4 --- Imperative: write a recursive function with a ref, from scratch
% =============================================================================

\quizproblem{Some Assembly Required}{5}{4}

Write \code{sumList : int list -> int}, which computes the sum of the
elements of a list, keeping a \emph{running total in a ref cell}.
Repeated calls must be correct.

\answerspace{5.5cm}

\begin{solution}
\begin{codeblock}
  fun sumList (L : int list) : int =
    let
      val total = ref 0
      fun go [] = !total
        | go (x::xs) = ( total := !total + x; go xs )
    in
      go L
    end
\end{codeblock}

\textbf{Rubric:} 1, one \code{ref} per top-level call; 2, recursion
accumulates via \code{:=} and the base case returns \code{!total}; 1,
correct overall.
\end{solution}

\newpage

% =============================================================================
% PROBLEM 5 --- Async: promises (upon timing + minimal-dependency rewrite)
% =============================================================================

\quizproblem{Promises, Promises}{8}{6}

A \term{promise} is a value representing the result of an asynchronous
operation that we do not yet have. Calling an asynchronous function
\emph{starts} its operation and \emph{immediately} returns a promise for
the result:

\begin{codeblock}
  val sleep     : int -> unit promise      (* n seconds     *)
  val read_line : unit -> string promise   (* line of input *)
\end{codeblock}

Promises are used through these operations, all of which also return
immediately:

\begin{itemize}
  \item \code{upon (p, k)} --- registers callback \code{k} to run once
    \code{p} resolves.
  \item \code{p thenDo f} --- once \code{p} resolves with \code{x}, runs
    \code{f x} (which returns another promise); the whole expression is a
    promise for that final result.
  \item \code{both (p, q)} --- a promise that resolves with the pair,
    once both \code{p} and \code{q} have resolved.
\end{itemize}

\textbf{(a)} \subpts{2} Which is printed first, \code{ping} or
\code{pong}? Justify in one sentence.

\begin{codeblock}
  ( upon (sleep 5, fn () => print "pong")
  ; print "ping"
  )
\end{codeblock}

\answerspace{2.5cm}

\begin{solution}
\textbf{\code{ping} prints first.} \code{upon} only registers the
callback and returns immediately, so \code{print "ping"} runs right away;
\code{pong} prints five seconds later, when the promise resolves.

\textbf{Rubric:} 1, \code{ping}; 1, \code{upon} returns immediately.
\end{solution}

\textbf{(b)} \subpts{4} This pipeline reads two stock tickers (e.g.
\code{"AAPL"}), fetches each one's current price with
\code{fetch_quote : string -> int promise}, then computes the spread
with \code{spread : int * int -> int promise}:

\begin{codeblock}
  read_line () thenDo (fn tick1 =>
  read_line () thenDo (fn tick2 =>
    both (fetch_quote tick1, fetch_quote tick2)
      thenDo (fn (p1, p2) => spread (p1, p2))))
\end{codeblock}

This computes the right answer, but some work waits on a result it does
not need. \textbf{Rewrite the pipeline so that every operation starts as
soon as the data it depends on is available}, still computing the spread
at the end.

\answerspace{5.5cm}

\begin{solution}
The false dependency: \code{fetch_quote tick1} needs only \code{tick1},
but it is not called until after the second ticker has been typed. Call
it as soon as \code{tick1} arrives:

\begin{codeblock}
  read_line () thenDo (fn tick1 =>
    both (fetch_quote tick1,
          read_line () thenDo (fn tick2 => fetch_quote tick2))
      thenDo (fn (p1, p2) => spread (p1, p2)))
\end{codeblock}
The first fetch now runs while the second ticker is being typed, and
\code{both} joins the two price promises. Also fine: bind \code{val q1 =
fetch_quote tick1}, then \code{read_line () thenDo (fn tick2 => both
(q1, fetch_quote tick2) thenDo ...)}. The call is what starts the work;
the promise just remembers it.

\textbf{Rubric:} 2, first fetch starts without waiting on the second
read (either shape); 1, second read-then-fetch via \code{thenDo}; 1,
\code{both} joins into \code{spread}.
\end{solution}

\end{document}
